% Generated draft: requires mathematical and source review.
\documentclass[11pt]{article}
\usepackage[margin=25mm]{geometry}
\usepackage{fontspec}
\setmainfont{Latin Modern Roman}
\usepackage{amsmath,amssymb,mathtools,unicode-math}
\setmathfont{Latin Modern Math}
\usepackage{enumitem,needspace,xurl,hyperref}
\setlength{\parindent}{0pt}
\setlength{\parskip}{5pt}
\setlength{\emergencystretch}{3em}
\newcommand{\N}{\mathbb{N}}
\newcommand{\Z}{\mathbb{Z}}
\newcommand{\Q}{\mathbb{Q}}
\newcommand{\R}{\mathbb{R}}
\newcommand{\C}{\mathbb{C}}
\newcommand{\F}{\mathbb{F}}
\newcommand{\A}{\mathbb{A}}
\newcommand{\PP}{\mathbb{P}}
\newcommand{\E}{\mathbb{E}}
\newcommand{\eps}{\varepsilon}
\newcommand{\dd}{\,\mathrm{d}}
\begin{document}
\section*{UnsolvedMath ID: 30003264}
\textbf{Source problem code:} catalogue entry
\subsection{Definitions and mathematical statement}
Let \(F\) be a global field, meaning a finite extension of \(\Q\) or of \(\F_p(t)\) for a prime \(p\). Let \(V_F\) be its set of places and \(S_\infty\) its set of archimedean places; thus \(S_\infty=\varnothing\) in the function-field case. For a finite set \(S\subseteq V_F\), put
\[
\mathcal O_S=\{a\in F: v(a)\geq 0\text{ for every nonarchimedean place }v\notin S\},
\]
where each nonarchimedean valuation is normalized to have value group \(\Z\). For such a place \(v\), let \(k(v)\) be the residue field of its valuation ring. Write \(R=\Z[1/2]\), and let \(H_3(G,R)\) denote group homology of the discrete group \(G\) with trivial coefficients in \(R\). The inclusion \(\mathcal O_S\hookrightarrow F\) induces a homomorphism
\[
\iota_{S,*}:H_3(\operatorname{SL}_2(\mathcal O_S),R)\longrightarrow H_3(\operatorname{SL}_2(F),R).
\]
Set \(K_S=\ker(\iota_{S,*})\).

The source uses \(\mathcal P(k)\) for a pre-Bloch group. A usual presentation is the abelian group generated by symbols \([x]\), for \(x\in k\setminus\{0,1\}\), subject to
\[
[x]-[y]+[y/x]
-\left[\frac{1-x^{-1}}{1-y^{-1}}\right]
+\left[\frac{1-x}{1-y}\right]=0
\qquad(x,y\in k\setminus\{0,1\},\ x\ne y).
\]
The supplied record does not specify whether this precise presentation, including its conventions for small residue fields, is intended. It also names a natural map \(\pi_S\) without giving its construction or a characterization determining it.

With those unresolved source qualifications, the supplied question is whether, for every global field \(F\), the map
\[
\pi_S:K_S\longrightarrow
\bigoplus_{v\in S\setminus S_\infty}
\bigl(\mathcal P(k(v))\otimes_{\Z}R\bigr)
\]
is an isomorphism for every sufficiently large finite set \(S\) of places. Interpreting ``sufficiently large'' as eventuality under inclusion, the intended quantifiers would be
\[
\forall F\ \exists S_0\subseteq V_F\text{ finite}\ \forall S\subseteq V_F\text{ finite},\qquad
S_0\subseteq S\ \Longrightarrow\ \pi_S\text{ is an isomorphism}.
\]
One may include \(S_\infty\) in \(S_0\) under this interpretation. The record does not explicitly specify this meaning of ``sufficiently large,'' any further restrictions on \(F\) or \(S\), or the definition of \(\pi_S\). Consequently, the displayed assertion remains a source question whose complete formulation requires clarification of these points.
\subsection{Short English statement}
For a global field \(F\), consider the third homology classes of \(\operatorname{SL}_2(\mathcal O_S)\), with coefficients in \(\Z[1/2]\), that become zero over \(F\). The question asks whether, once the finite set of places \(S\) is sufficiently large, the source's natural map identifies this kernel with the direct sum of residue-field pre-Bloch groups at the nonarchimedean places in \(S\), after tensoring with \(\Z[1/2]\). The map, precise pre-Bloch conventions, and meaning of sufficiently large require source clarification.
\subsection{Sources}
\begin{itemize}
\item Supplied local UnsolvedMath registry record, permanent ID \(30003264\), exact source problem code \texttt{catalogue entry}, title: ``Third Homology and Pre-Bloch Groups of \(S\)-Arithmetic \(\operatorname{SL}_2\).'' Collection: Oberwolfach Reports Open Problems. No individual proposer is supplied; the proposed year is recorded as 2016.
\item Supplied source citation: \emph{Mini-Workshop: Computations in the Cohomology of Arithmetic Groups} (2017), Oberwolfach Reports, pp.\ 2941--2973, DOI \texttt{10.4171/owr/2016/52}. Source URL: \url{https://doi.org/10.4171/owr/2016/52}. The supplied background attributes the problem to workshop open-problem discussions documented in Oberwolfach Reports.
\item Supplied literature reference, labeled ``OWR 52 report'': \url{https://doi.org/10.4171/owr/2016/52}. Literature checked at: \texttt{2026-08-21}. This is supplied dated curation, not an independently verified assessment. It labels the problem open and reports that no authoritative later source located in that curation proves or disproves eventual isomorphy of \(\pi_S\). The background explicitly requests independent verification of time-sensitive status and cited claims.
\item The registry labels the statement \texttt{corrected\_verified} and says its mathematical content was preserved while notation, typography, or omitted definitions were repaired. This is a supplied editorial qualification, not independent mathematical verification of this draft.
\end{itemize}
\end{document}
